%========================================================================================
%  Seção 16.3 do D2L -- Sentiment Analysis: Using Convolutional Neural Networks
%  https://d2l.ai/chapter_natural-language-processing-applications/sentiment-analysis-cnn.html
%========================================================================================

%------------------------------------------------
\section{Análise de sentimento com CNN}
%------------------------------------------------

\begin{frame}{A matemática da convolução 1D}
    \begin{columns}[c]
        \begin{column}{0.44\textwidth}
            \footnotesize
            Uma \destaque{lupa de $k$ palavras} percorre a frase: em cada parada, ela multiplica o que vê pelo molde e soma.

            \vspace{0.15cm}
            \begin{itemize}\setlength{\itemsep}{3pt}
                \item \textbf{Input}: o texto já é número --- cada célula, uma palavra.
                \item \textbf{Kernel}: o \destaque{molde} de $k$ números que procura um n-grama.
                \item \textbf{Output}: um número por parada --- o encaixe daquele pedaço.
                \item Os números do kernel são \textbf{aprendidos}, não escolhidos.
            \end{itemize}
        \end{column}
        \begin{column}{0.54\textwidth}
            \centering
            \begin{tikzpicture}[font=\small,
                cel/.style={draw=dexpiBlue, fill=dexpiBlueLight, minimum width=0.85cm, minimum height=0.50cm, inner sep=0pt},
                hl/.style={draw=dexpiGreenDark, fill=dexpiGreenLight, minimum width=0.85cm, minimum height=0.50cm, inner sep=0pt, font=\small\bfseries},
                kn/.style={draw=dexpiGreenDark, fill=dexpiGreenLight, minimum width=0.85cm, minimum height=0.50cm, inner sep=0pt, font=\small\bfseries},
                lbl/.style={anchor=east, font=\scriptsize\bfseries, text=dexpiBlue}]
                \node[hl] (in0) at (0.425,0) {0};
                \node[hl] (in1) at (1.275,0) {1};
                \node[cel] (in2) at (2.125,0) {2};
                \node[cel] (in3) at (2.975,0) {3};
                \node[cel] (in4) at (3.825,0) {4};
                \node[cel] (in5) at (4.675,0) {5};
                \node[cel] (in6) at (5.525,0) {6};
                \node[lbl] at (-0.12,0) {Input};

                \node[kn] (k0) at (0.425,-1.05) {1};
                \node[kn] (k1) at (1.275,-1.05) {2};
                \node[lbl] at (-0.12,-1.05) {Kernel};
                \draw[-{Stealth}, thick, dexpiRedDark] (2.10,-1.05) -- (3.20,-1.05);
                \node[font=\scriptsize, text=dexpiRedDark, anchor=south] at (2.65,-1.05) {desliza};

                \node[hl] (out0) at (0.85,-2.10) {2};
                \node[cel] at (1.70,-2.10) {5};
                \node[cel] at (2.55,-2.10) {8};
                \node[cel] at (3.40,-2.10) {11};
                \node[cel] at (4.25,-2.10) {14};
                \node[cel] at (5.10,-2.10) {17};
                \node[lbl] at (-0.12,-2.10) {Output};

                \node[draw=dexpiBlue, fill=dexpiBlueLight, rounded corners=2pt, inner sep=4pt, font=\small] (eq) at (1.70,-3.25)
                    {${\color{dexpiBlue}0}\cdot{\color{dexpiGreenDark}1} + {\color{dexpiBlue}1}\cdot{\color{dexpiGreenDark}2} = {\color{dexpiGreenDark}\mathbf{2}}$};
                \draw[-{Stealth}, thick, dexpiGreenDark] (eq.north) -- (out0.south);
            \end{tikzpicture}

            \vspace{0.1cm}
            {\scriptsize No texto real, cada célula é um vetor de 100 números e o kernel tem 100 números por palavra; aqui usamos escalares para ver a conta.\par}
        \end{column}
    \end{columns}

    \note{
        Pensem numa lupa de duas palavras que percorre a frase. O texto já é número: 0, 1, 2, e assim por diante. Em cada parada, a lupa multiplica o que vê pelos números do molde e soma. Aqui: 0 vezes 1 mais 1 vezes 2 dá 2, o primeiro número do output. Depois ela desliza uma casa e repete. O detalhe importante: os números do molde não são escolhidos por nós --- a rede os aprende, e eles passam a disparar quando o padrão aparece.
    }
\end{frame}

\begin{frame}{A arquitetura textCNN na prática}
    \begin{columns}[c]
        \begin{column}{0.52\textwidth}
            \footnotesize
            A mesma conta, agora dentro de um modelo de verdade \cite{kim2014convolutional}.

            \vspace{0.15cm}
            \begin{itemize}\setlength{\itemsep}{3pt}
                \item A RNN lê \textbf{palavra por palavra}: o estado de hoje carrega o de ontem.
                \item A CNN olha \textbf{pedaços} de 3, 4 e 5 palavras \textbf{ao mesmo tempo}.
                \item \textbf{Fluxo}: GloVe $\to$ CNN $\to$ positivo/negativo.
                \item \textbf{Expressões locais}: ``não é bom'' é um trigrama --- a CNN o detecta como padrão.
            \end{itemize}
        \end{column}
        \begin{column}{0.44\textwidth}
            \centering
            \begin{tikzpicture}[font=\scriptsize,
                bx/.style={draw=dexpiBlue, fill=dexpiBlueLight, rounded corners=2pt, align=center, minimum width=1.85cm, minimum height=1.35cm, inner sep=2pt},
                mid/.style={draw=dexpiBlue, fill=dexpiBlue, text=white, rounded corners=2pt, align=center, minimum width=1.85cm, minimum height=1.35cm, inner sep=2pt, font=\scriptsize\bfseries}]
                \node[bx] (g) at (0,0) {Pré-treino\\{\tiny GloVe (15.5)}\\{\tiny 6 bi tokens}};
                \node[mid] (m) at (2.4,0) {Arquitetura\\{\tiny CNN (16.3)}\\{\tiny conv + max}};
                \node[bx] (a) at (4.8,0) {Aplicação\\{\tiny sentimento}\\{\tiny pos./neg.}};
                \draw[-{Stealth}, line width=1.5pt, dexpiRedDark] (g) -- (m);
                \draw[-{Stealth}, thick, dexpiBlue] (m) -- (a);
                \node[font=\tiny, text=dexpiRedDark, align=center, anchor=south] at (1.20,0.85) {aqui entra a\\convolução};
                \draw[dexpiRedDark, dashed, thick] (1.20,0.78) -- (1.20,0.06);
            \end{tikzpicture}
        \end{column}
    \end{columns}

    \vspace{0.1cm}
    \begin{exampleblock}{Por que é mais rápida}
        \footnotesize
        Todas as janelas são calculadas \textbf{em paralelo}; a RNN não pode, porque cada passo espera o anterior.
    \end{exampleblock}

    \note{
        A RNN lê palavra por palavra: o estado de hoje depende do de ontem. A CNN não: ela olha pedaços de três, quatro e cinco palavras ao mesmo tempo, e todas as janelas são calculadas em paralelo --- é por isso que ela é mais rápida. O fluxo é o desta figura: o GloVe entrega os vetores prontos, a CNN os transforma em escores, e a camada final classifica. E onde ela brilha é no local: ``não é bom'' é um trigrama, e a CNN o detecta como um padrão.
    }
\end{frame}

\begin{frame}{Max-over-time: muitos números, um resumo}
    \begin{columns}[c]
        \begin{column}{0.50\textwidth}
            \footnotesize
            Depois da convolução, cada filtro deixa uma \textbf{sequência} de escores. Como resumir a frase?

            \vspace{0.15cm}
            \begin{itemize}\setlength{\itemsep}{1pt}
                \item Para cada canal, guardamos \textbf{só o maior} escore: um número por filtro.
                \item Onde ele apareceu não importa: o filtro de ``não é bom'' dá 8 no início, no meio ou no fim.
                \item É um \destaque{filtro de importância}: o pico é a evidência de que o padrão existe.
                \item A média \textbf{diluiria} --- uma pista que aparece uma vez em 500 posições vale $1/500$; daí o \textit{global max pooling} de Collobert et al. \cite{collobert2011natural}.
            \end{itemize}
        \end{column}
        \begin{column}{0.48\textwidth}
            \centering
            \begin{tikzpicture}[font=\scriptsize]
                \foreach \v [count=\i from 0] in {1,5,2,8,3} {
                    \ifnum\i=3
                        \node[draw=dexpiGreenDark, fill=dexpiGreenLight, minimum width=0.72cm, minimum height=0.6cm, font=\scriptsize\bfseries] (v\i) at (0.95*\i,0) {\v};
                    \else
                        \node[draw=dexpiBlue, fill=dexpiBlueLight, minimum width=0.72cm, minimum height=0.6cm] (v\i) at (0.95*\i,0) {\v};
                    \fi
                }
                \node[font=\tiny, text=dexpiBlue, anchor=south] at (1.9,0.34) {escores de um filtro, um por posição};
                \draw[-{Stealth}, thick, dexpiGreenDark] (v3.south) -- (2.85,-0.62);
                \node[draw=dexpiGreenDark, fill=dexpiGreenLight, rounded corners=2pt, font=\scriptsize, anchor=north] at (2.85,-0.62) {max $=8$};
                \node[font=\tiny, text=gray, anchor=north] at (2.85,-1.18) {média $=3{,}8$};
            \end{tikzpicture}
        \end{column}
    \end{columns}

    \vspace{0.1cm}
    \begin{exampleblock}{O que o pooling compra}
        \footnotesize
        Saída de tamanho \textbf{fixo}: $C$ números, um por canal, para qualquer $T$.
    \end{exampleblock}

    \note{
        A convolução deixou uma sequência de escores por filtro: são muitos números. Para resumir, guardamos só o maior de cada canal. Reparem que a posição não importa: se o filtro procura ``não é bom'', o valor 8 diz ``achei em algum lugar'', seja no início, no meio ou no fim. É um filtro de importância. E por que não a média? Porque ela dilui: uma pista que aparece uma vez em quinhentas posições vale um quinhentos do total. O pico preserva o sinal.
    }
\end{frame}

\begin{frame}{A arquitetura completa do textCNN}
    \begin{columns}[c]
        \begin{column}{0.50\textwidth}
            \footnotesize
            O modelo de Kim (2014), de baixo para cima.

            \vspace{0.15cm}
            \begin{itemize}\setlength{\itemsep}{3pt}
                \item \textbf{Entrada}: a frase vira matriz --- 11 palavras $\times$ 6 dimensões no livro.
                \item \textbf{Filtros de larguras diferentes}: 2 e 4 no livro (3, 4 e 5 no nosso código), cada um atrás de um n-grama.
                \item \textbf{Pooling $+$ concatenação}: cada filtro vira 1 número; $4+5=9$ no exemplo.
            \end{itemize}
        \end{column}
        \begin{column}{0.48\textwidth}
            \centering
            \begin{tikzpicture}[font=\tiny]
                \fill[dexpiBlueLight] (0,0) rectangle (2.2,1.2);
                \draw[dexpiBlue!55, step=0.2] (0,0) grid (2.2,1.2);
                \node[below, font=\tiny] at (1.1,0) {11 palavras $\times$ 6};

                \fill[dexpiGreenLight] (0,1.9) rectangle (2.0,2.7);
                \draw[dexpiGreenDark!60, step=0.2] (0,1.9) grid (2.0,2.7);
                \fill[dexpiGreenLight] (2.4,1.7) rectangle (4.0,2.7);
                \draw[dexpiGreenDark!60, step=0.2] (2.4,1.7) grid (4.0,2.7);
                \node[anchor=east, font=\tiny, text=dexpiGreenDark] at (-0.05,2.3) {$k=2$};
                \node[anchor=west, font=\tiny, text=dexpiGreenDark] at (4.05,2.3) {$k=4$};

                \draw[-{Stealth}, thick, dexpiBlue] (1.0,1.2) -- (1.0,1.9);
                \draw[-{Stealth}, thick, dexpiBlue] (2.0,1.2) -- (3.0,1.7);

                \draw[-{Stealth}, dashed, thick, dexpiRedDark] (1.0,2.7) -- (1.0,3.1);
                \draw[-{Stealth}, dashed, thick, dexpiRedDark] (3.2,2.7) -- (3.2,3.1);
                \node[draw=dexpiRedDark, fill=dexpiRedLight, font=\scriptsize, minimum width=0.5cm] at (1.0,3.35) {4};
                \node[draw=dexpiRedDark, fill=dexpiRedLight, font=\scriptsize, minimum width=0.5cm] at (3.2,3.35) {5};
                \node[anchor=west, font=\tiny, text=dexpiRedDark] at (1.32,3.35) {max};
                \node[anchor=west, font=\tiny, text=dexpiRedDark] at (3.52,3.35) {max};

                \node[draw=dexpiBlue, fill=dexpiBlueLight, rounded corners=2pt, font=\tiny, inner sep=3pt] (c) at (2.1,4.0) {concat $=9$ $\rightarrow$ densa $\rightarrow$ 2};
                \draw[-{Stealth}, thick, dexpiBlue] (1.0,3.62) -- (1.55,3.82);
                \draw[-{Stealth}, thick, dexpiBlue] (3.2,3.62) -- (2.65,3.82);
            \end{tikzpicture}
        \end{column}
    \end{columns}

    \note{
        Vamos guiar o olhar de baixo para cima. Embaixo, a frase vira uma matriz: onze palavras por seis dimensões no exemplo do livro. Em cima dela, duas convoluções trabalham em paralelo: uma com filtros de duas palavras, outra com filtros de quatro. São tamanhos diferentes de propósito --- cada uma captura um n-grama diferente. As setas pontilhadas sobem para o pooling: cada canal vira um número. No fim, quatro mais cinco dão nove, que é o vetor que vai para a classificação.
    }
\end{frame}

\begin{frame}{O fluxo de dados: de $T$ palavras a 2 números}
    \begin{columns}[c]
        \begin{column}{0.48\textwidth}
            \footnotesize
            Os tamanhos, sem susto.

            \vspace{0.15cm}
            \begin{itemize}\setlength{\itemsep}{3pt}
                \item Convolução de largura $k$: a saída tem $T-k+1$ posições ($11-2+1=10$; $11-4+1=8$).
                \item O pooling reduz cada canal a \textbf{1} número; a concatenação junta todos ($4+5=9$; no nosso, $3 \times 100=300$).
                \item \textbf{Dropout} desliga neurônios ao acaso no treino; a camada densa mapeia o vetor para \textbf{2 logits}.
            \end{itemize}
        \end{column}
        \begin{column}{0.50\textwidth}
            \centering
            \begin{tikzpicture}[font=\tiny]
                \foreach \i in {0,...,8} {
                    \ifnum\i<4
                        \fill[dexpiGreenLight] ({0.30*\i},0) rectangle ++(0.27,0.27);
                        \draw[dexpiGreenDark] ({0.30*\i},0) rectangle ++(0.27,0.27);
                    \else
                        \fill[dexpiBlueLight] ({0.30*\i},0) rectangle ++(0.27,0.27);
                        \draw[dexpiBlue] ({0.30*\i},0) rectangle ++(0.27,0.27);
                    \fi
                }
                \draw[decorate, decoration={brace, amplitude=3pt, mirror}] (0,-0.05) -- (1.2,-0.05) node[midway, below=4pt] {4};
                \draw[decorate, decoration={brace, amplitude=3pt, mirror}] (1.2,-0.05) -- (2.7,-0.05) node[midway, below=4pt] {5};
                \node[anchor=south, font=\tiny, text=dexpiBlue] at (1.35,0.35) {vetor concatenado};

                \draw[-{Stealth}, thick, dexpiBlue] (2.75,0.13) -- (3.15,0.13);
                \node[draw=dexpiRedDark, dashed, fill=white, font=\tiny, minimum width=0.95cm, minimum height=0.5cm, inner sep=1pt] at (3.7,0.13) {dropout};
                \draw[-{Stealth}, thick, dexpiBlue] (4.2,0.13) -- (4.55,0.13);
                \node[draw=dexpiBlue, fill=dexpiBlueLight, font=\tiny, minimum width=0.75cm, minimum height=0.5cm, inner sep=1pt] at (4.95,0.13) {densa};
                \draw[-{Stealth}, thick, dexpiBlue] (5.35,0.13) -- (5.7,0.13);
                \node[anchor=west, font=\tiny] at (5.75,0.13) {2 logits};
            \end{tikzpicture}

            \vspace{0.15cm}
            {\scriptsize
            \begin{tabular}{@{}ll@{}}
                \toprule
                etapa & forma \\
                \midrule
                entrada & $11 \times 6$ \\
                conv $k=2$ / $k=4$ & $10 \times 4$ / $8 \times 5$ \\
                pooling & $4+5=9$ \\
                densa $+$ dropout & $2$ \\
                \bottomrule
            \end{tabular}\par}
        \end{column}
    \end{columns}

    \note{
        Agora os tamanhos, sem susto. Uma convolução de largura k sobre T palavras devolve T menos k mais um posições: onze com filtro dois dá dez; com filtro quatro, oito. O pooling reduz cada canal a um número só --- quatro mais cinco, nove. É esse vetor que entra na camada densa, passando antes pelo dropout, que desliga neurônios ao acaso no treino para o modelo não decorar. A saída são dois logits, positivo e negativo.
    }
\end{frame}

\begin{frame}[fragile]{O código do textCNN}
\begin{columns}[c]
    \begin{column}{0.62\textwidth}
\begin{lstlisting}[basicstyle=\ttfamily\fontsize{6.1}{7.1}\selectfont,aboveskip=0pt,belowskip=0pt]
class TextCNN(nn.Module):
    def __init__(self, vocab, d, ks, cs):
        super().__init__()
        self.emb = nn.Embedding(vocab, d)        # ajustavel
        self.fix = nn.Embedding(vocab, d)        # congelado
        self.convs = nn.ModuleList(
            [nn.Conv1d(2 * d, c, k) for c, k in zip(cs, ks)])
        self.pool = nn.AdaptiveMaxPool1d(1)      # max no tempo
        self.drop = nn.Dropout(0.5)
        self.out = nn.Linear(sum(cs), 2)

    def forward(self, inputs):                   # (B, T)
        E = torch.cat((self.emb(inputs), self.fix(inputs)), dim=2)
        E = E.permute(0, 2, 1)                   # (B, 200, T)
        H = [self.pool(torch.relu(c(E))).squeeze(-1)
             for c in self.convs]                # 3 x (B, 100)
        return self.out(self.drop(torch.cat(H, dim=1)))  # (B, 2)

d, ks, cs = 100, [3, 4, 5], [100, 100, 100]
\end{lstlisting}
    \end{column}
    \begin{column}{0.34\textwidth}
        \footnotesize
        \textbf{Onde está cada ideia}
        \begin{enumerate}\setlength{\itemsep}{2pt}
            \item \texttt{emb} $+$ \texttt{fix}: os 200 canais.
            \item \texttt{Conv1d}: um detector por canal e por largura.
            \item \texttt{pool}: o máximo no tempo, um número por canal.
            \item \texttt{cat} $+$ \texttt{out}: 300 números $\to$ 2 logits.
        \end{enumerate}
        \vspace{0.15cm}
        {\scriptsize\color{dexpiRedDark} GloVe: \texttt{copy\_} nas duas tabelas e \texttt{requires\_grad = False} em \texttt{fix}.}
    \end{column}
\end{columns}

\note{
    Este é o código, com os nomes encurtados. No construtor, as duas tabelas, a lista de convoluções, o pooling, o dropout e a camada de saída. No forward: concatenamos as duas tabelas, transpomos, aplicamos cada convolução com ReLU e pooling, e concatenamos os resultados.
    Uma observação para quem for comparar com o livro: o texto chama o pooling de max-over-time, mas o listing em PyTorch do D2L traz AdaptiveAvgPool1d, que faz média. Aqui usamos o máximo, que é o que o texto descreve; se você reproduzir o código do livro como está, vai treinar com média.
    Os hiperparâmetros são os do livro: embedding de 100, larguras 3, 4 e 5, 100 canais, dropout 0,5, Adam com taxa 10 menos 3, 5 épocas.
}
\end{frame}

\begin{frame}{Resultados: CNN $\times$ BiRNN, e os limites}
    \begin{columns}[c]
        \begin{column}{0.46\textwidth}
            \centering
            \footnotesize
            \begin{tabular}{lcc}
                \toprule
                 & BiRNN & textCNN \\
                \midrule
                perda de treino & $\approx 0{,}49$ & $0{,}066$ \\
                acurácia treino & $\approx 0{,}80$ & $0{,}979$ \\
                acurácia teste & $\approx 0{,}83$ & $\mathbf{0{,}868}$ \\
                \bottomrule
            \end{tabular}
            \vspace{0.15cm}
            {\scriptsize IMDb; GloVe, $d = 100$; CNN: Adam, $\text{lr}=10^{-3}$, 5 épocas.}
        \end{column}
        \begin{column}{0.50\textwidth}
            \footnotesize
            \begin{itemize}\setlength{\itemsep}{2pt}
                \item O teste sobe de $\approx 0{,}83$ para $\mathbf{0{,}868}$, com os mesmos dados e o mesmo GloVe; mudam a arquitetura e a taxa ($10^{-3}$ contra $10^{-2}$).
                \item \textbf{O alarme}: treino $0{,}98$ contra teste $0{,}87$. A CNN decora mais --- \texttt{dropout} $0{,}5$ e metade dos embeddings congelada são as defesas.
                \item \textbf{Por que funciona}: o \textit{max} ignora \textbf{onde} a pista está; ``is so great'' vale no começo ou no fim.
            \end{itemize}
        \end{column}
    \end{columns}

    \vspace{0.1cm}
    \begin{alertblock}{O limite da janela}
        \footnotesize
        Um filtro de 3 palavras não liga ``apesar do começo lento'' a ``é ótimo'' no fim de uma crítica longa. Por isso Kim usa 3, 4 e 5; a evolução natural é a \textbf{atenção} (BERT, Seção 15.8), que liga qualquer posição a qualquer outra.
    \end{alertblock}

    \note{
        Agora a comparação, com os números do livro. A BiRNN chega a cerca de 83 por cento de acurácia no teste, com perda de treino de 0,49. O textCNN chega a 86,8, com perda de 0,066. Os mesmos dados e os mesmos embeddings; mudam a arquitetura e a taxa de aprendizado, 10 menos 3 no lugar de 10 menos 2.
        Mas olhem a distância entre treino e teste: 0,98 contra 0,87. A CNN decora boa parte do treino. As duas defesas estão no modelo: o dropout de 0,5 e a metade congelada dos embeddings. E vale lembrar que a acurácia de treino da RNN era medida ao longo da época, então a comparação de gap não é perfeitamente limpa; o que importa é o teste.
        O limite é estrutural: um filtro de 3 palavras só vê 3 palavras. Ele não liga "apesar do começo lento" a "é ótimo" no fim de uma crítica longa. A solução de Kim é ter filtros de 3, 4 e 5; a solução seguinte, e a que domina hoje, é a atenção: cada posição olha todas as outras, que é exatamente o BERT que vimos na seção 15.8.
    }
\end{frame}
